\subsection{Correlazione}\label{sec:results_correlation}


%L'analisi Within Country ha mostrato come, per i Paesi che hanno ospitato le
%Olimpiadi nelle edizioni dal 1964 al 2024, ospitare l'evento abbia portato in
%tutti i casi tranne uno (USA 1996) un numero maggiore di medaglie rispetto alla
%loro performance media.

L'analisi Within Country dal 1964 al 2024 ha mostrato come ospitare abbia in
tutti i casi portato significativamente più medaglie (in percentuale, tra +0.9 e
+3.4, o +50.2\% e +373.4\%), tranne che per USA (-1.3, o -9.9\%) e Canada (solo
+0.1, o +3.7\%).


% il testo basta a descrivere i dati, non consumerei due slot tabella per questo
% e spearman. L'ho messo in result.

% comunque, unire la germania con quella ovest non è tanto corretto, come ho
% scritto da altre parti (sono confini diversi, Paesi diversi...)


% info generale da conclusione

%Tuttavia, analizzando l'andamento del medagliere per alcuni di questi Paesi, si
%può notare come nelle edizioni precedenti rispetto a quella ospitata ci sia
%stato un incremento graduale.

% informazioni generali, le diremo insieme per tutti

%Il picco avuto nell'edizione ospitata potrebbe avere una spiegazione nei
%maggiori investimenti nello sport fatti negli anni precedenti, anche in vista
%della notizia dell'accettazione da parte del comitato Olimpico della
%candidatura da parte di quello specifico Paese. Infatti, ai Paesi viene
%comunicato che saranno loro ad ospitare le Olimpiadi per una specifica edizione
%ben 7 anni prima, dando allo Stato tempo per iniziare ad investire nelle
%infrastrutture sportive.


% sposto l'interpretazione qui, così, stando vicino, si capiscono i risultati

%Il Test dei Segni ha confermato quanto può essere visionato anche ad occhi: il
%test ha calcolato infatti un p-value di 0.0034

Il \textbf{Test dei Segni} conferma il risultato: il \textbf{$p$-value di
0.0034} (sintomo di sbilanciamento tra i segni) permette di rigettare con
estrema confidenza l'ipotesi nulla $H_0$ e confermare l'esistenza di un
vantaggio sistematico nell'ospitare.


%\subsection*{Correlazione PIL/Popolazione Medaglie}


% in questo studio ci poniamo domande, non tentiamo di dimostrare niente in
% particolare (nel senso di far valere un'ipotesi).

%I coefficienti di Spearman calcolati hanno in parte confermato quanto tentiamo
%di dimostrare in questo studio: infatti, sia la correlazione popolazione $->$
%medaglie, sia la correlazione PIL $->$ medaglie ha mostrato valori positivi.

% frase da conclusione non serve anticipare.

%Questi numeri ci mostrano come ricchezza e bacino di talenti spieghino
%abbastanza il successo di un Paese alle Olimpiadi. In seguito, grazie all’uso
%di una regressione multipla andremo a verificare che l'``effetto host'' non ha
%un grande potere sul successo di un Paese alle Olimpiadi.
%
%Abbiamo calcolato i coefficienti di Spearman annuali: \begin{itemize} \item la
%correlazione popolazione $->$ medaglie mostra valori che variano dallo 0.43
%allo 0.54. Tali valori possono essere interpretati come una correlazione
%moderata \item la correlazione PIL $->$ medaglie mostra valori che variano
%dallo 0.64 allo 0.69. Tali valori mostrano una correlazione forte \end{itemize}
%
%\subsubsection*{Correlazione PIL $->$ percentuale di medaglie vinte}
%\begin{verbatim} Year 1964 0.657628 1968 0.664404 1972 0.688908 1976 0.676049
%1988 0.644296 1992 0.686920 1996 0.652784 2000 0.670057 2004 0.683186 2008
%0.673528 2012 0.668123 2016 0.685533 2020 0.651082 2024 0.649676 \end{verbatim}
%
%\subsubsection*{Correlazione Popolazione $->$ percentuale di medaglie vinte}
%\begin{verbatim} Year 1964 0.492175 1968 0.504393 1972 0.541326 1976 0.506511
%1988 0.503505 1992 0.492770 1996 0.514203 2000 0.483528 2004 0.509060 2008
%0.500601 2012 0.434991 2016 0.480924 2020 0.467844 2024 0.431100 \end{verbatim}



I coefficienti di Spearman (che variano da -1 a 1) hanno mostrato una certa
correlazione in entrambe le applicazioni: moderata per popolazione $->$ medaglie
(che varia dallo 0.43 allo 0.54); forte per PIL $->$ medaglie (tra 0.64 e 0.69).

